\documentclass[10pt]{article}
\usepackage[utf8]{inputenc}
\usepackage[T1]{fontenc}
\usepackage{amsmath}
\usepackage{amsfonts}
\usepackage{amssymb}
\usepackage[version=4]{mhchem}
\usepackage{stmaryrd}
\usepackage{graphicx}
\usepackage[export]{adjustbox}
\graphicspath{ {./images/} }

\begin{document}
если $X$ - сепарабельное гильбертово пространотво, $x \in X, y \in X,\left\{e_{n}\right\}$ - ортонормированный базис в $X$, $a_{n}=\left(x, e_{n}\right)$ и $b_{n}=\left(y, e_{n}\right)$ - коэффициенты Фурье соответственно әлементов $x$ и $y$, то справедливо равенство

\[
(x, y)=\sum_{n=1}^{\infty} a_{n} \bar{b}_{n}
\]

наз. о б о бщ енны м р а в енс т в о П а р с ев а л я. В достаточно законченном віде вопрос о полноте систем функций, являющихся собственными функциями дифференциальных операторов, был изучөн В. А. Стекловым [1].

П. р. обобщцается и на случай несепарабельных гильбертовых пространств: если $\left\{e_{\alpha}\right\}, \alpha \in \mathfrak{A}(\mathfrak{A}$ нек-рое множество индексов), является полной ортонормированной системой гильбертова пространства $X$, то для любого элемента $x \in X$ справедливо П. $\mathrm{p}$.

\[
(x, x)=\sum_{\alpha \in \mathfrak{A}}\left|\left(x, e_{\alpha}\right)\right|^{2},
\]

причем сумма в правой части равенства понимается как

\[
\sup _{\mathfrak{P}_{0}} \sum_{\alpha \in \mathfrak{A}}\left|\left(x, e_{\alpha}\right)\right|^{2},
\]

где верхняя грань берется по всевозможным конечным подмножествам $\mathfrak{U}_{0}$ множества $\mathfrak{Q}$.

В случае, когда $X=L_{2}[-\pi, \pi]$ состоит из действительных функций, квадрат к-рых интегрируем по Лебегу на отрезке $[-\pi, \pi], f \in L_{2}[-\pi, \pi]$, в качестве полной ортогональной системы взята тригонометрич. система функций и

\[
f \sim \frac{a_{0}}{2}+\sum_{n=1}^{\infty} a_{n} \cos n x+b_{n} \sin n x,
\]

равенство (1) имеет вид

\[
\frac{1}{\pi} \int_{-\pi}^{\pi} f^{2}(t) d t=\frac{a_{0}^{2}}{2}+\sum_{n=1}^{\infty}\left(a_{n}^{2}+b_{n}^{2}\right)
\]

и наз. кл а с си ч е с ким р а вен с т вом П а рс е в а л я; оно было указано М. Парсевалем (M. Parseval, 1805).

Если $g \in L_{2}[-\pi, \pi]$ и

\[
g \sim \frac{a_{0}^{\prime}}{2}+\sum_{n=1}^{\infty} a_{n}^{\prime} \cos n x+b_{n}^{\prime} \sin n x,
\]

то равенство, аналогичное формуле (2), выглядит следуюпим образом:

\[
\frac{1}{\pi} \int_{-\pi}^{\pi} f(t) g(t) d t=\frac{1}{2} a_{0} a_{0}^{\prime}+\sum_{n=1}^{\infty}\left(a_{n} a_{n}^{\prime}+b_{n} b_{n}^{\prime}\right) .
\]

Классы $K$ и $K^{\prime}$ действительных функций, определенных на отрезке $[-\pi, \pi]$, такие, что для всех $f \in K$ и $g \in K$ имеет место обобщенное П. р. (3), наз. дополнительными. Примером дополнительных классов являются пространства $L_{p}[-\pi, \pi]$ и

\[
L_{p^{\prime}}[-\pi, \pi], \frac{1}{p}+\frac{1}{p^{\prime}}=1, \quad 1<p<+\infty .
\]

Лum.: [1] С т е н л о в В. А., «Записки физико-математич. общества», сер. 8,1904, т. 15 , № 7 , с. 1-32; [2] Н и к о л ьс к и й С. М., Курс математического анализа, 2 изд., т. 2 , М.,

\includegraphics[max width=\textwidth, center]{2023_07_08_967a5a8b95a356728587g-1}
Тригонометрические ряды, М., 1961; [5] 3 и г м у н д А., Тригонометрические ряды, пер. с англ., т. 1, M., $1965 ;$ [6] К и р и л ло в А. А., Г в иш и а н и А. Д., 'Теоремы и за-
дачи функционального анализа, М., 1979.

ПАРСЕВАЛЯ - ПЛАНШЕРЕЛЯ ФОРМУ ЛА - сМ. Планшереля теорема.

ПАСКАЛЕВА ГЕОМЕТРИЯ - геометрия плоскости, построенной над полем (коммутативным телом). Название этой геометрии связано с тем, что в этой геометрии на плоскости выполняется конфигурационное предло же ние П а п п а точки $1,3,5$ п 2, 4, 6 соответственно лежат на прямых (коллинеарны), то точки пересечения пар прямых $(1,2)$ и $(4,5),(2,3)$ и $(5,6),(3,4)$ и $(6,1)$ - точки 9,7 , 8 - также лежат на одной прямой при любом выборе системы образующих точек $1,3,5$ на од-
ной прямой и 2,4, 6 - на другой прямой, отличной от первой (см. рис.). Важнейший частный случай этого предложения (в афффинной плоскости) утверждает: из того,

\includegraphics[max width=\textwidth, center]{2023_07_08_967a5a8b95a356728587g-1(1)}
что прямая $(4,3)$ параллельна $(6,5)$ и прямая $(6,1)$ параллельна $(2,3)$, следует параллельность прямых $(2,5)$ и $(4,1)$.

П. г. плоскости может быть построена над бесконечными или конечными полями, соответственно этому плоскость наз. бесконечной или конечной п а с к ал е в о й п ло с к о с т ь ю.

Впервые важную роль предложения Паскаля в построении геометрич. систем над бесконечными полями исследовал Д. Гильберт (D. Hilbert, см. [1]), к-рый устанавливал доказуемость предложения Паскаля при различных наборах аксиом из системы аксиом евклидова пространства. Д. Гильберт показал, что Паскаля теорему в бесконечной плоскости можно доказать с помощью плоскостных аксиом инцидентности, порядка, конгруәнтности, параллельности и непрерывности, причем было установлено, что без аксиом непрерывности в этом случае теорему Паскаля доказать нельзя.

Опираясь же на пространственные аксиомы системы Гильберта, теорему Паскаля можно доказать без аксиом конгруэнтности, но обязательно с применением аксиом непрерывности (исключение аксиом непрерывности в бесконечной плоскости приводит к непаскалевой геометрии). Возможность доказательства теоремы Паскаля аналогична в указанном смысле возможности доказательства Дезарга предложения с использованием пространственных аксиом, однако в доказательстве теоремы Паскаля проявляется особая роль аксиом непрерывности Архимеда в бесконечных плоскостях (см. Неархимедова геометрия).

Предложение Паппа - Паскаля проективно выполняется в нек-рой плоскости тогда и только тогда, когда умножение во всех натуральных телах этой плоскости обладает коммутативным свойством, или иначе: натуральное тело всякой паскалевой плоскости является полем и, наоборот, плоскость, построенная над полем, обладает П. г. Поэтому П. г. иногда наз. г е оме т р и е й с комм у т а тивным умно жен и е м. Таким образом, в паскалевой плоскости конфигурационное предложение Паппа - Паскаля имеет алгебраич. иввариант, выражающий коммутативное свойство умножения в множестве, над к-рым построена эта плоскость.

В любой проективной плоскости теорема Паппа Паскаля влечет за собой предложение Дезарга.

Конечная паскалева плоскость как конечная проективная плоскость существует только в случае, если число точек, лежащих на каждой прямой этой плоскости, есть $p^{s}+1$, где $p$ - простое, $s$ - натуральное число. Так как всякое конечное альтернативное тело является полем, то в конечной плоскости предложение Дезарга влечет за собой теорему Паппа - Паскаля, причем последняя является следствием т. н. малой теоремы Дезарга. Вместе с тем существуют конечные проективные плоскости, являющиеся непаскалевыми. Паскалева плоскость пзоморфна двойственной себе плоскости.


\end{document}